% chapters/ch03.tex —— 第 3 章
% 体例：具体失败 → 设计问题 → 机制解释 → 对照实验 → 条件变化 → 选择原则
% 数字身份纪律：
%   - E01 运行 E01-20260923T152230Z-90acda 与任务集 v1 的计数一律 \measured{}；
%   - 2602.19320（Anatomy of Agentic Memory）、2607.12227（Rethinking Harness Evolution）一律 \cited{}；
%   - 由上述数字算出的倍数、差值、比值一律 \calcnum{}。
% 本章是全书的方法论枢纽：后面的对照全部依赖本章的协议。任何判据的放宽都须进 agent/decision-record.md。
\chapter{如何检验持续改进}
\label{ch03}

\chaptergoal{一个系统第二次做得比第一次好，成绩提高来自更新机制，还是比较条件变了？}
\chapterprops{P1关注预算在持久更新与任务内重试之间怎样分配，P1b关注继续更新相对只读复用的增量。本章用四格试验说明对照、成本、区间与结果解释的关系；P2所问的因素贡献另需统一交叉设计。}
\chapterfig{图~\ref{fig:ch03-budget}：同一份预算的四种花法，以及各自的收益留在哪一层}
\chapterpitfall{只展示最好的一次运行——更隐蔽的版本是：展示一组挑过题目的“平均提升”，
这个均值可能不代表原定任务分布}

% ---------- 1. 具体失败 ----------
\begin{failurebox}
在一批 \measured{16} 道历史封存题上，我们让同一个 Agent 在四个对手格里每题重复运行 \measured{5} 次：
四格使用同一模型、题集、执行工具和回合上限，
分别比较静态执行、外部只读经验、任务专属适配与跨任务学习的经验。
相同上限并不保证完整总预算相等；两类记忆的来源也不同，后文将逐项审计。

结果读起来很像一次成功：
静态执行格的回合通过率 \measured{0.7375}、
只读的持久基线 \measured{0.7250}、
每道题单独适配 \measured{0.7625}、
跨任务持久更新 \measured{0.7875}。
最后一格最高，相对第一格的点估计是 \measured{+0.0500}——五个百分点。

\textbf{点估计可以如实报告，但“持久更新更好”的推断现在还不能成立。}
同一组数据给出的 95\% 自助区间是 \measured{$[-0.1000, +0.1750]$}：它跨过 0。
在格4的这 \measured{16} 道题里，有 \measured{9} 道题在 \measured{5} 次重复运行中出现了
通过 / 未通过不一致。单次结果存在运行波动，不能仅凭一次成功判断某题已经获得稳定改善。

问题同时出在估计精度与比较设计。
按该次试验的配对差标准差 \measured{0.2875}，以三个百分点为目标区间半宽，
历史精度公式给出 \calcnum{353} 道题的规划近似；它不是已经注册的功效保证。
当前只有 \measured{16} 道题，而且完整重试基线和总成本归属仍需核验。
\end{failurebox}

\designquestion{自进化系统的评测有一个需要特别处理的困难：被比较的对象常常不是两个独立模型，
而是\textbf{同一个系统在两个时刻}——前一次的运行会改变后一次的起点。
于是“第二次更好”至少有四种解释：真的学到了、题目被挑过、题目恰好更简单、
或者只是多试了一次。本章的问题是：
\textbf{怎样安排一次比较，使成绩变化可以被归属到一个机制上，并且这个归属可以被推翻。}}

\begin{sloppypar}
第 2 章给了三层时间尺度与 $A_k$ 这套写法，回答的是“改了什么”。
本章回答后半句：“凭什么认为它有用”。
它为后续机制比较提供四项阅读线索——
被固定住的比较条件（\S\ref{sec:ch03-control}）、四类题集（\S\ref{sec:ch03-sets}）、
四项指标（\S\ref{sec:ch03-metrics}）、配对比较与区间（\S\ref{sec:ch03-paired}），
再用四类对照检查收益应归属到哪个方案（\S\ref{sec:ch03-fourarms}）。历史试验的落实程度需要逐项核对，不能因本章写出了要求，就认定每张结果表均已满足。
\end{sloppypar}

% ---------- 2. 控制什么 ----------
\section{控制什么：四件必须同时固定的东西}
\label{sec:ch03-control}

\begin{sloppypar}
比较两个版本的系统时，至少要同时固定四件事。
它们不是“实验设计的最佳实践清单”，而是四条独立的混淆来源：漏掉任何一条，
观测到的差都可能与更新机制无关。
\end{sloppypar}

\begin{table}[htbp]
\centering
\caption{四条控制项：不固定它会混淆什么、在本书实验台落在哪个字段}
\label{tab:ch03-control}
\begin{tabular}{>{\raggedright\arraybackslash}p{2.4cm}>{\raggedright\arraybackslash}p{4.2cm}>{\raggedright\arraybackslash}p{3.8cm}>{\raggedright\arraybackslash}p{3.0cm}}
\toprule
控制项 & 不固定会混淆什么 & 本书的落点 & 违反时的后果 \\
\midrule
基础模型 & 能力差被读成机制差 & \texttt{config.json} 的 \texttt{model} / \texttt{base\_url} & 无法单独归因于更新机制 \\
\midrule
工具与执行环境 & 工具能力差被读成学到的技能 & 核对实际工具清单、依赖与隔离边界；持久写入字段只检查已登记路径 & 工具差异被读成学到的技能 \\
\midrule
任务分布 & 题目难度差被读成迁移能力 & 任务集 v1：\measured{8} 族 $\times$ \measured{8} 题 = \measured{64} 题，划分只按族内编号 & 挑选引入分布偏差 \\
\midrule
预算 & 多花的算力被读成更好的机制 & 主账与逐回合轨迹记录调用、步数、token和计时；步数保存在轨迹中，计时范围见下文 & “预算更多”冒充“方法更好” \\
\bottomrule
\end{tabular}
\end{table}

实现的检查粒度也属于预算口径。现有执行器在下一轮模型调用前检查累计用量，因此一次调用可能使token或墙钟读数越过阈值；它没有在远程请求中途按剩余额度实时截断。调用上限、记账阈值与服务端最大生成长度应分别检查。表中的预算约束不能被读成所有物理资源都已具有严格的实时硬上限，详细实现限制见附录\ref{app:lab}。


\begin{sloppypar}
四条里边最容易被跳过的是第一条：换主干模型等于换实验。
这一控制项有具体的外部例子：一项记忆系统分析的总体比较包含六种系统，但其主干敏感性表只报告其中两种\cite{anatomy-memory-2602}。
将\texttt{gpt-4o-mini}换成\texttt{Qwen-2.5-3B}后，
同一个系统的答题分从 \cited{0.289} 掉到 \cited{0.102}（\texttt{SimpleMem}），
另一个从 \cited{0.781} 掉到 \cited{0.447}（\texttt{Nemori}），
降幅并不是同一个比例。这里的主干同时参与答题和记忆更新、整合，读写条件一起变化。
\end{sloppypar}

\begin{sloppypar}
原作者将结构化记忆操作的不稳定解释为一种可能不在短期对话中显露的失效机制。
表4的格式列具体记录\emph{可恢复的格式偏差}，数值更高表示更依赖后备解析；它不是不可恢复记忆损坏的直接发生率。
SimpleMem的该比例从\cited{1.20\%}升到\cited{4.82\%}，Nemori从\cited{17.91\%}升到\cited{30.38\%}。
答题分与格式偏差一同变化支持检查接口适配的重要性，但仅凭这张表，不能把全部答题下降定量归因于写入损坏。
这条实证给了本书一条可以直接执行的纪律：
\textbf{任何跨模型的收益比较，必须先跑该模型的静态基线}。
第 5 章把它做成了 E03c 的一整个臂（\texttt{base\_cross}），原因就在这里。
\end{sloppypar}

\begin{principlebox}
控制项的作用是让“只有一个变量不同”这句话可以被检查。
写一个控制项时，要说明它如何记录、凭什么证据核查，以及偏离时怎样处理。
能由程序强制的约束应有相应检查，依赖外部环境或人工核查的条件应保留版本与审计记录。
\end{principlebox}

\subsection{先写清楚究竟比较什么}

E01 的格 4 得到更高的题均通过率，并不能独自回答“记忆有没有用”。
若参照物是完全没有记忆的系统，比较同时包含记忆内容、检索方式和构建过程；
若参照物是同一份记忆的只读版本，比较才更接近“继续更新是否值得”。
读者必须先写出参照物，再解释差值，否则一句“提升”会在几个不同问题之间滑动。

\textbf{待估计量}是实验准备估计的那个具体差异。
设目标任务分布为 $\mathcal D$，总资源约束为 $b$，
比较两个完整方案 $A$ 与 $B$，本章关注的量可以写成
\begin{equation}
  \Delta(\mathcal D,b)
  =\mathbb E_{x\sim\mathcal D}\left[q(A,x;b)-q(B,x;b)\right].
  \label{eq:ch03-estimand}
\end{equation}
若执行包含采样，式中的$q$表示给定任务时对声明范围内的运行随机性取平均的质量。逐题重跑用于估计给定版本的条件均值；若研究经验构建、更新和候选选择本身的波动，则应重复相应完整流程。方案包含经验来源、更新过程、候选选择与最终执行，不能只留下测试时的那一次模型调用。
这个式子定义了要问的问题，本身不保证现有题集能够代表 $\mathcal D$。

给一个 Agent 人工编写的经验库，另一个 Agent 从训练失败中自己积累经验，
测到的差异包括“谁提供了经验”。
若想隔离写入机制，应另设共同起始记忆，并让两个方案从同一个版本分叉。
现有 E01 格 2 使用外部固定条目，格 4 使用训练中产生的条目；
因此格 4 减格 2 是两个完整方案之差，不能直接命名为纯粹的写入收益。
对应实现见 \texttt{e01\_arms.py} 的 \texttt{build\_external\_memory} 与
\texttt{run\_e01.py} 的训练阶段。

观察范围也必须进入问题。
任务族内新题上的效果、未见任务族上的效果、换读取者后的效果各有独立的目标分布。
在一张表里同时列出这些数字可以帮助读者比较，但把它们平均成一个“迁移率”会抹掉检验对象。
第 5 章的两套拆分正是为不同问题准备的，不能因为都叫测试题就交换使用。

\subsection{相同预算上限不等于相同总投入}

两臂都允许相同次数的调用，实际花费仍可能不同。
一臂早早完成，另一臂把所有调用用尽；
一臂发送短提示，另一臂每次附带长记忆；
一臂复用既有工件，另一臂在测试前还生成了提示或筛选了技能。
相同调用上限控制了一个资源维度，没有控制全部资源。

总成本应按发生位置分解，而不是按是否处于用户等待时间内分解：
\begin{equation}
  C_{\mathrm{total}}=
  C_{\mathrm{source}}+C_{\mathrm{update}}+C_{\mathrm{select}}
  +C_{\mathrm{test}}+C_{\mathrm{maintain}}.
  \label{eq:ch03-cost-scope}
\end{equation}
五项分别对应来源采集、更新、选择、测试与维护；失败尝试归入发生它的那一项。
没有可靠价格时，每项保留调用数、提示与输出 token、墙钟时间，
不能将不同物理量随意相加为一个无单位总成本。

评审反馈也要计量资源。Kumar等在成对偏好评审中，给定候选对、评审模型与呈现顺序，Direct与G-Eval每次判断各调用API\cited{1}次，FLASK分解评审调用\cited{4}次。成对评审的Direct输出上限为\cited{32} token，四种JSON/XML格式均为\cited{4096} token。调用数与允许的输出长度共同界定了这些流程的资源条件\cite{all-verdicts-2610-aacl}。

本书需要区分两张图。
固定测试预算的质量图回答：已经获得这个版本以后，它执行任务的效果如何。
计入更新过程的总预算图回答：从相同起点出发，花同样资源最终能得到什么效果。
前一张图允许讨论成熟版本的部署价值，后一张图才直接服务 P1 的预算分配主张。
两张图可以同时报告，图题必须说明成本范围。

共享来源轨迹的生成成本也不能凭空消失。
在同一批经验上比较两种表示形式时，可以另列一次共同成本，再分别列各臂的增量成本；
若比较两个独立部署方案，则各方案均须承担它依赖的来源成本。
摊销规则在运行前声明，不能看见哪种记账方式有利于新方法后再选择。
历史 E01 已有主账与副账，但完整的总预算对等仍需要单独审计，不能由相同回合上限直接推出。


% ---------- 3. 四类题集 ----------
\section{四类题集：训练、选择验证、回归、封存}
\label{sec:ch03-sets}

\begin{sloppypar}
本书区分四种数据用途，并在任务集v1中分别划分题目。其他设计未必也要建立四个物理目录，但需要明确哪些数据参与更新和选择，哪些数据支持最终评价：
\end{sloppypar}

\begin{table}[htbp]
\centering
\caption{四类题集：用途、纪律与本书任务集 v1 的对应划分}
\label{tab:ch03-sets}
\begin{tabular}{>{\raggedright\arraybackslash}p{2.6cm}>{\raggedright\arraybackslash}p{4.0cm}>{\raggedright\arraybackslash}p{3.6cm}>{\raggedright\arraybackslash}p{2.8cm}}
\toprule
题集 & 用途 & 硬纪律 & v1 划分 \\
\midrule
训练任务 & 产生经验，形成候选 & 预先规定来源与划分，不按最终结果挑选 & \texttt{train}，\measured{16} \\
\midrule
选择验证任务 & 在候选之间做选择 & 预设查询上限；换数据须另存新划分 & \texttt{select-valid}，\measured{16} \\
\midrule
旧能力回归集 & 检测退化 & 预先规定旧能力范围，固定前后比较条件 & \texttt{regression}，\measured{16} \\
\midrule
封存测试集 & 报告收益 & 最终评测启动前，开发与选择阶段的查询次数必须为 0 & \texttt{sealed}，\measured{16} \\
\bottomrule
\end{tabular}
\end{table}

\begin{sloppypar}
任务集 v1 有 \measured{8} 个族、每族 \measured{8} 题，共 \measured{64} 题。
族的定义是“一类失效模式”（如“可变默认参数与共享可变状态”），不是“一类题目”；
族内同构是刻意的——第 5 章要用“在 A 族上积累的经验迁到 B 族还剩多少”来区分两种拟合目标。
划分只按族内编号（\texttt{taskset\_spec.py} 的 \texttt{SPLIT\_BY\_INDEX}），
本次划分不按模型的最终成绩挑题。若事后只留下有利题目，目标分布与报告分母就会改变；固定编号能避免这种特定选择，却不能证明生成的题目已代表实际部署任务。
\end{sloppypar}

\begin{sloppypar}
查询预算 $Q$ 是选择验证集那一栏的核心。
本书把它实现成一个会抛异常的对象（\texttt{eval.py} 的 \texttt{QueryBudget}）：
该对象能阻止超过登记额度的查询，但\texttt{rotate}只重置计数器，不会自动生成新任务。真正轮换需要另提供未参与开发的数据并保存划分版本。选择验证集本来就参与候选选择，反复使用增加过拟合风险；它不能再兼任独立最终评测。
E01四格没有接入该查询计数器；后续实验必须分别核对调用链，不能因公共模块存在就宣称查询隔离已经实现。
\end{sloppypar}

% ---------- 4. 四项指标 ----------
\section{四项指标：成功率、成本、迁移、遗忘}
\label{sec:ch03-metrics}

\begin{sloppypar}
只报通过率的评测会漏掉三件坏事：成本高得不可接受、换了入口就失效、旧能力悄悄退化。
本书要求每次更新同时报四项。
\end{sloppypar}

\begin{table}[htbp]
\centering
\caption{四项指标：问什么、怎么算、只看它会漏掉什么}
\label{tab:ch03-metrics}
\begin{tabular}{>{\raggedright\arraybackslash}p{2.0cm}>{\raggedright\arraybackslash}p{4.2cm}>{\raggedright\arraybackslash}p{3.8cm}>{\raggedright\arraybackslash}p{3.4cm}}
\toprule
指标 & 问的是 & 怎么算 & 只看它会漏掉什么 \\
\midrule
成功率 & 新题做对了吗 & 封存集 pass@1，题为单位取均值 & 其余三项 \\
\midrule
成本 & 变好值这个价吗 & 调用次数 / token / 墙钟；副账（构建持久物）单列 & 用十倍算力换三个点 \\
\midrule
迁移 & 换个入口还在吗 & 换模型或换 harness 后的 $\Delta_{\text{migrate}}$ & 记忆只对当前模型有效 \\
\midrule
遗忘 & 旧能力还在吗 & 旧任务历史最佳减当前的均值 $FGT_t$ & 提升是拆东墙补西墙 \\
\bottomrule
\end{tabular}
\end{table}

\begin{sloppypar}
成本这一项常被写成脚注，但它能改变结论。
同一份实证\cite{anatomy-memory-2602}把用户可见时延拆成检索与生成两段并列表给出（表 5）：
直塞全上下文的时延为 \cited{1.726} 秒，而一个带层级结构的系统在同等设置下达到 \cited{32.372} 秒，
约为前者的 \calcnum{18.8} 倍；同一批系统里也有 \cited{1.057} 秒的。
也就是说，“加了记忆”这个动作的系统开销可以相差一个数量级以上，
而通过率的表里看不出这件事。
\end{sloppypar}

\begin{sloppypar}
还有一类更隐蔽的成本：维护。
该表量化了用户可见时延，也另列离线构建时间与token；\textbf{维护成本则被省略}，
因为这部分常常被异步处理掉。
这类费用应另列账目，并在报告中说明覆盖。历史E01把任务适配与经验写入的副调用同时并入主账的调用数和token，\texttt{side\_*}是其中分项，不能再与主账相加。副调用未单独记录墙钟，回合计时也未覆盖这些任务外等待；因此该副账仍不代表持续维护费。
异步不等于免费，缺少计量也不等于费用为零。
\end{sloppypar}

\subsection{终点成功与过程可执行性}
\label{sec:ch03-prefix-executability}

成功率还需要过程诊断。Prefix Executability 的作者构造了一个通讯录修复示例：Agent 把电话号码当作联系人标识请求删除，参数虽然也是字符串，当前状态中却没有这个标识，动作被拒绝；随后查到正确标识再删除，任务仍能完成\cite{prefix-executability-2610}。这是原文附录C.2的示意例，不是本书实测。最终成功记录了修复后的结果，首次失效记录了另一件事；两者都应保留。

这里的\textbf{前缀可执行性}检查的是：沿此前动作形成的状态，每一步是否满足接口、数据依赖与状态前提。格式能够解析，只完成了其中一层检查。原文将按基准状态语义执行的确定性评分，与执行前由语言模型检查计划的验证器分开；计划获准不代替实际轨迹验收。它也不独自认证任务语义、安全或隐私，部分可观测环境还需要额外的状态估计\cite{prefix-executability-2610}。

读这类曲线时尤其要核对分母。设第$i$条轨迹有$T_i$个被评价的动作位置，首次失效位置为$F_i$；全程可执行时规定$F_i=T_i+1$。原文固定批次的经验曲线为
\begin{equation}
 \widehat{\mathrm{PEC}}(k)=\frac{1}{N}\sum_{i=1}^{N}
 \mathbf{1}\{T_i\ge k,\ F_i>k\}.
 \label{eq:ch03-prefix-executability}
\end{equation}
分母$N$保留全部被评价轨迹。因此曲线下降既可能来自失效，也可能来自轨迹已经结束；它不是纯粹的错误率。若改用到达该深度的轨迹数$n_k=\sum_i\mathbf{1}\{T_i\ge k\}$作分母，得到的是另一个条件问题，$n_k=0$时该比例没有定义\cite{prefix-executability-2610}。

用本书构造的三个轨迹说明：一条在第二步后仍有效，一条第一步就合法完成，另一条到达第二步却在那里失效。深度为二时，固定批次比例是\calcnum{1/3}，只在到达该深度的轨迹中计算则是\calcnum{1/2}。前者保留起始批次，后者改变了条件集合；不能把两条曲线的差当成方法优劣。给本书实验增加这类诊断，还需保存足以还原状态与工具接受条件的轨迹，不能仅凭最终布尔判定反推。

% ---------- 5. 配对比较 ----------
\section{配对比较、重复运行与不确定性}
\label{sec:ch03-paired}

\begin{sloppypar}
本书的默认比较方式是\textbf{配对差}：同一道题上跑两个版本，取差，再在题上平均。
它让两个版本在相同任务条件下形成对应比较；方差能减少多少取决于两臂结果的相关性。
\end{sloppypar}

\begin{enumerate}
  \item \textbf{重复运行}。同一配置跑 $R$ 次，取均值。$R$ 不是形式主义的重复：
        实测显示在这 \measured{16} 道题上，四个格里有 \measured{9} 到 \measured{12} 道题
        在 \measured{5} 次运行中出现通过 / 未通过不一致。单次结果可以作为个案记录，若要估计该题上的随机成功概率，则需要重复与不确定性说明。
  \item \textbf{重采样单位}。默认以题为单位自助；若怀疑族内相关，改以族为单位。
        重采样单位应同时服从抽样结构与推断目标；族内相关不能因目标叫“新题”就忽略。只有很少几个族时，换成族级重采样也无法补出不存在的覆盖。
  \item \textbf{区间而非裸判定}。报点估计加 95\% 自助区间。
        本书要求同时说明差异大小与估计精度，区间用于展示待估计差异的不确定范围。
  \item \textbf{不挑选运行}。$R$ 次全部进账，不许删掉“网络不好那次”。
\end{enumerate}

\subsection{重复运行降低哪一种不确定性}

同一道题多跑几次，能更清楚地看见该题的成功概率；它不会自动增加任务覆盖。
设第 $i$ 道题、第 $r$ 次运行的二值结果为 $Y_{ir}$，
每题的通过率为 $\bar Y_i=R^{-1}\sum_rY_{ir}$。
比较两臂时先得到题内差 $d_i=\bar Y_{i,A}-\bar Y_{i,B}$，
再在题之间求均值，不能把同一题的重复运行伪装成多道独立新题。

在“各题独立抽取、给定题目后的重复运行独立同分布”的教学模型下，
记题内成功概率为 $p_i$，题均通过率的方差可分解为
\begin{equation}
 \operatorname{Var}(\bar Y_i)
 =\operatorname{Var}(p_i)+\frac{\mathbb E[p_i(1-p_i)]}{R}.
 \label{eq:ch03-repeat-variance}
\end{equation}
这个式子来自全方差分解：第一项描述题与题之间的差异，第二项描述同一题重复调用的波动。
增加 $R$ 只压低第二项；要改善第一项的估计，仍需增加符合目标分布的独立题目。
真实服务若共享缓存、随时间改变模型或存在调用相关性，上述独立性假设必须重新检查。

配对也不是自动消掉全部难度差。
由 $\operatorname{Var}(X-Y)=\operatorname{Var}(X)+\operatorname{Var}(Y)-2\operatorname{Cov}(X,Y)$ 可知，
两臂在同一题上正相关时，配对往往能减少比较方差；若相关性很弱或为负，收益就不同。
NIST 对配对观察的定义和差值统计给出了一般形式\cite{nist-paired}；
本书在此基础上把每题的重复运行均值作为配对对象。
运行器记录的种子还须与后端实现核对：种子被写入日志，并不意味着服务端按该种子确定性采样。

\subsection{留出任务族以后，仍然需要足够多的族}

E03 拆分 A 留出了两个任务族，确实使经验来源和测试任务在族身份上分离。
但这只证明测试题来自这两个未参与训练的族，
不能证明它们足以代表所有未来任务族。
族内题目共享构造方式和失效结构，许多相似变体可能提供相近的信息。

统计单位应在看结果前与推广范围一起确定。
若目标是这批固定族中的新题，可以报告条件于这些族的题级结果，并明确限制；
若目标是未见过的族，就需要多个独立留出族，以及与族级抽样相符的分析。
把重采样选项从题改成族，并不能凭空补出新的族；只有两个族时，族级不确定性的估计仍十分脆弱。

任务集扩容因此有两种不同动作。
增加同族变体主要改善已有失效模式内的估计，增加新的独立族才扩大机制覆盖。
计划书应分别列出族数、每族题数、每题重复次数，以及对各族的权重。
一个族因容易生成而拥有更多题时，直接题均会让它获得更高权重；
如果工程目标要求每族等权，应在协议中明确采用族均再平均。

\subsection{区间跨零、工程阈值与反证不是同一件事}

E01 的正向点估计和跨零区间可以同时成立。
它们分别描述样本中心和估计的不确定范围，没有逻辑冲突。
跨零意味着当前区间同时包含两种方向，不自动给出“没有效果”，
也不能仅凭它断言唯一问题必定是样本少：真实效应接近零、方差较大或设计不匹配也可能造成这种结果。

工程阈值回答的是“多大的改善值得付出成本”，统计区间回答的是“数据把差异限制到什么范围”。
本书保留原先三个百分点的设计目标，不因当前结果难判而调高或调低。
点估计超过该目标但区间很宽，仍然不能证明真实改善已经超过目标。
若整个有效区间都低于一个预先声明的最低有用改善，可以反对相应的工程主张；
这也不等于证明精确的零效应。

声称两种方法在工程上等价，需要先定义可接受差异的范围，并按相应的等价性设计收集数据。
“区间跨零”和“区间完整落入可接受范围”是两种不同关系。
前者可以出现在一个非常宽、兼容明显正负效应的区间里，后者要求足够精确的估计。
该次试验未预先冻结等价界限与检验，因此不能事后选一个足够宽的范围给 P3 判等价。

\texttt{conditional} 也不能充当所有不确定结果的收纳箱。
它要求已经有证据支持一个具体条件边界，例如在预先规定的某种反馈下得到支持，
在另一种反馈下存在可信的不同表现，并且差异本身得到合适检验。
“一组显著、另一组不显著”并不自动证明两组效应不同；
第 4 章讨论反馈强度时，需要直接比较相应效应，而不是只比较两行状态标签。

\subsection{区间精度与检验功效要分开规划}

旧版用 $n\ge(1.96s/\mathrm{MDE})^2$ 估算题数，再把它称为“判决力”。
这个称呼遗漏了漏检风险。
在独立样本、近似正态且标准差可用的条件下，
该式估算的是双侧区间半宽不超过给定值 $h$ 所需的规模；
NIST 将这种精度规划与同时给定显著性水平和功效的设计分开列出\cite{nist-sample-size}。

检验功效是在指定真实效应下，检验拒绝相应零假设的概率。
若规划双侧检验的显著性水平 $\alpha$、目标功效 $1-\beta$，
并把希望检出的真实均值差记为 $\delta_{\mathrm{detect}}$，
一个常用的正态近似为
\begin{equation}
 n_{\mathrm{power}}\approx
 \left(\frac{(z_{1-\alpha/2}+z_{1-\beta})s}
 {\delta_{\mathrm{detect}}}\right)^2.
 \label{eq:ch03-power}
\end{equation}
此式用于围绕零差异的均值检验；若检验目标是超过非零工程门槛，分母应反映目标真值与零假设边界的距离。
它还没有处理小样本标准差不确定性、任务族相关、多重比较或序贯观察。

以双侧 $\alpha=0.05$、目标功效 $0.80$、检出差异 $0.03$ 作为\textbf{规划示例}，
代入 E01 已记录的未舍入标准差，格 4 减格 1 约需 \calcnum{721} 道独立配对题，
格 4 减格 2 约需 \calcnum{1708} 道。
这是计算结果，不是新增实验，也不是已经批准的正式样本量。
复算入口为 \texttt{calculations/audit\_sample\_size.py}，它同时保留原精度估算，
使读者能看见增加功效要求后规模为何改变。

先用试验数据估方差，再为下一批数据冻结规模，是合理的工作顺序。
把这一规模追溯写成产生试验数据之前已经注册，则改变了事实。
正式方案还应为族间差异、预计缺失、比较数量及运行成本留出设计空间，
并说明何时可以停止；不能每天看一次区间，直到它第一次不跨零才结束。


% ---------- 6. 泄漏与过拟合 ----------
\section{评测泄漏、选择偏差与验证集过拟合}
\label{sec:ch03-leak}

\begin{sloppypar}
持久更新增加了一条需要核查的信息通路：系统可能把见过的评价材料写下来带到下一次。是否比其他系统更容易泄漏仍需比较，不能从结构直接推断发生率。
三类问题按发生顺序排列如下。
\end{sloppypar}

\begin{table}[htbp]
\centering
\caption{三类评测失真与设计防线；具体实现还须核对调用链}
\label{tab:ch03-leak}
\begin{tabular}{>{\raggedright\arraybackslash}p{2.6cm}>{\raggedright\arraybackslash}p{4.4cm}>{\raggedright\arraybackslash}p{4.0cm}>{\raggedright\arraybackslash}p{2.4cm}}
\toprule
类型 & 发生在哪一步 & 典型症状 & 本书的防线 \\
\midrule
评测泄漏 & 检验题进入了更新过程 & 更新后的系统见过这些题 & 开发阶段不查询封存题；隔离仍待落实 \\
\midrule
选择偏差 & 同一批题既用于选候选又用于报收益 & 提升随查询次数一起涨 & 选择与终评分开；查询计数须实际接入 \\
\midrule
验证集过拟合 & 反复在同一批题上挑 & 曲线单调上升，换一批题归零 & 到上限停止；新数据身份另行核对 \\
\bottomrule
\end{tabular}
\end{table}

\begin{sloppypar}
\textbf{基准饱和}影响的是改进空间，并不等同于信息泄漏。
若全部证据都能装进当前上下文，就需要把全上下文方案纳入比较，才能判断额外记忆结构改善了什么。
记忆系统分析论文据此建议报告同一主干上的\cite{anatomy-memory-2602}
\emph{上下文饱和差}（原文记作 $\Delta$，定义为增强记忆的系统与硬塞全上下文的基线在同一协议下的分数差），
并明确说它是一条\textbf{诊断量而不是及格线}。
该文Table 2按信息规模、交互深度、实体多样性估计常用基准的结构性饱和风险，
其中上下文只有千余 token 的一类被判为“对上下文窗口而言是平凡的”，
而另一个约 \cited{100000} token、能整体装进 \cited{128000} 上下文窗口的基准
被判为“高风险”。这些是基于结构统计的启发式判断，不是已测得满分或证明外部记忆无用；即使质量差很小，效率、更新便利与鲁棒性仍可另行比较。
\end{sloppypar}

\begin{sloppypar}
最后一类失真出在打分这一端。
该文把词汇类指标与语义裁判放在六个系统上对照（表 3），发现两者给出的排序不一致：
一个偏抽象摘要的系统在 F1 下排第五（\cited{0.116}），在语义裁判下排第四；
另一个 F1 拿到 \cited{0.268} 的系统，语义分数不到 \cited{0.30}。
作者把这类失配归为两种：\emph{Paraphrase Penalty}（正确的抽象回答因为字面重合低而被扣分）
与 \emph{Negation Trap}（高重合掩盖了事实错误）。
\end{sloppypar}

语义评审也要连同评分规则来读。表3中，Nemori与MAGMA的领先次序随评审提示互换。附录D.3给出的规则包含分级评分与二值判定，报告应将提示、评分尺度与排序一并列明，便于核对分数差究竟在比较什么\cite{anatomy-memory-2602}。

\begin{sloppypar}
对本书的直接后果是：对于本书可执行编程题，优先采用行为验收，不能仅用词面重合度替代；
任务集 v1 的历史运行由 \texttt{task.verify} 执行验收命令，
用例的期望值由构建期那个正确实现算出，不手写。
\end{sloppypar}

% ---------- 7. 四臂判据 ----------
\section{收益归属的四臂判据}
\label{sec:ch03-fourarms}

\begin{sloppypar}
即使控制、题集、指标、统计四件事都做对了，还剩最后一个混淆：
\textbf{成绩提高可能只是“多试了几次”或“多改了几轮”，与持久更新无关}。
四类对照按持久范围组织，帮助定位信息与更新的来源（图~\ref{fig:ch03-budget}）。严格的预算分配比较还需对齐共同起点、候选选择和完整资源范围；历史E01只落实了其中一部分。
\end{sloppypar}

% figures/ch03-budget.tex —— 第 3 章核心图（全书引用入口）
% 用法：在 book/chapters/ch03.tex 中 % figures/ch03-budget.tex —— 第 3 章核心图（全书引用入口）
% 用法：在 book/chapters/ch03.tex 中 % figures/ch03-budget.tex —— 第 3 章核心图（全书引用入口）
% 用法：在 book/chapters/ch03.tex 中 \input{figures/ch03-budget}
% 图主体在 ch03-budget-body.tex，便于单独编译与复用。
\begin{figure}[htbp]
\centering
\input{figures/ch03-budget-body}
\caption{历史E01四格的状态保留范围。静态重复取均值不同于生成多个候选后选择；固定记忆也可能帮助未见任务。四格分类不自动证明持久更新优越，完整预算与共同起点仍须核对。}
\label{fig:ch03-budget}
\end{figure}

% 图主体在 ch03-budget-body.tex，便于单独编译与复用。
\begin{figure}[htbp]
\centering
% figures/ch03-budget-body.tex —— 第 3 章核心图主体：历史四格的状态保留范围
% 图主体仅供 figures/ch03-budget.tex 引用；数字不进图，因此不受身份纪律约束。
% 全书引用入口：figures/ch03-budget.tex（含 figure 环境与编号）
% 图内不含任何数字，因此不受 \cited{} / \calcnum{} / \measured{} 身份纪律约束。
% 画的是第 3 章的收益归属判据：分类状态保留范围；收益需要独立对照。
\begin{tikzpicture}[
  blk/.style={draw, rounded corners=2pt, minimum height=9mm, text width=25mm, align=center, font=\footnotesize},
  bud/.style={blk, fill=gray!14},
  epi/.style={blk, fill=orange!12},
  task/.style={blk, fill=yellow!12},
  ver/.style={blk, fill=blue!12},
  keep/.style={blk, fill=green!12},
  drop/.style={blk, fill=red!8},
  lbl/.style={font=\footnotesize, align=left, text width=52mm},
  arrow/.style={-{Stealth[length=1.6mm]}, thick}
]
  % ---- 预算 ----
  \node[bud] (b) at (0,0) {比较协议\\预算单位、准备范围与信息来源};

  % ---- 四种花法 ----
  \node[epi]  (g1) at (3.4, 3.2) {历史格 1\\重复运行，报告平均};
  \node[epi]  (g2) at (3.4, 1.1) {历史格 2\\预先给定，此后只读};
  \node[task] (g3) at (3.4,-1.1) {历史格 3\\任务内适配，用完即弃};
  \node[ver]  (g4) at (3.4,-3.2) {历史格 4\\跨任务持久记忆};

  \draw[arrow] (b) -- (g1);
  \draw[arrow] (b) -- (g2);
  \draw[arrow] (b) -- (g3);
  \draw[arrow] (b) -- (g4);

  % ---- 收益留到哪 ----
  \node[drop] (d1) at (6.8, 3.2) {不新增跨题状态};
  \node[drop] (d2) at (6.8, 1.1) {保留固定记忆};
  \node[drop] (d3) at (6.8,-1.1) {当前任务状态};
  \node[keep] (d4) at (6.8,-3.2) {保留到后续版本};

  \draw[arrow] (g1) -- (d1);
  \draw[arrow] (g2) -- (d2);
  \draw[arrow] (g3) -- (d3);
  \draw[arrow] (g4) -- (d4);

  % ---- 判据 ----
  \node[lbl, text width=37mm] at (10.5, 3.2) {提供静态参照；\\尚无重试候选选择};
  \node[lbl, text width=37mm] at (10.5, 1.1) {检查固定信息作用；\\需核对是否同源};
  \node[lbl, text width=37mm] at (10.5,-1.1) {检查任务内适配；\\不能据此声称跨题收益};
  \node[lbl, text width=37mm] at (10.5,-3.2) {检查持久方案；\\是否有收益仍需比较};

  % ---- 底部 ----
  \draw[gray!45] (-1.5,-4.6) -- (12.5,-4.6);
  \node[lbl, text width=140mm] at (5.5,-5.3) {
    设计要求：\textbf{（一）}预算单位与全流程成本范围明确；
    \textbf{（二）}各臂依赖的准备成本均入账；
    \textbf{（三）}任务内状态与跨题持久物分开；
    \textbf{（四）}开发改善不替代独立最终评价。分类图本身不证明预算对齐或收益。};
\end{tikzpicture}

\caption{历史E01四格的状态保留范围。静态重复取均值不同于生成多个候选后选择；固定记忆也可能帮助未见任务。四格分类不自动证明持久更新优越，完整预算与共同起点仍须核对。}
\label{fig:ch03-budget}
\end{figure}

% 图主体在 ch03-budget-body.tex，便于单独编译与复用。
\begin{figure}[htbp]
\centering
% figures/ch03-budget-body.tex —— 第 3 章核心图主体：历史四格的状态保留范围
% 图主体仅供 figures/ch03-budget.tex 引用；数字不进图，因此不受身份纪律约束。
% 全书引用入口：figures/ch03-budget.tex（含 figure 环境与编号）
% 图内不含任何数字，因此不受 \cited{} / \calcnum{} / \measured{} 身份纪律约束。
% 画的是第 3 章的收益归属判据：分类状态保留范围；收益需要独立对照。
\begin{tikzpicture}[
  blk/.style={draw, rounded corners=2pt, minimum height=9mm, text width=25mm, align=center, font=\footnotesize},
  bud/.style={blk, fill=gray!14},
  epi/.style={blk, fill=orange!12},
  task/.style={blk, fill=yellow!12},
  ver/.style={blk, fill=blue!12},
  keep/.style={blk, fill=green!12},
  drop/.style={blk, fill=red!8},
  lbl/.style={font=\footnotesize, align=left, text width=52mm},
  arrow/.style={-{Stealth[length=1.6mm]}, thick}
]
  % ---- 预算 ----
  \node[bud] (b) at (0,0) {比较协议\\预算单位、准备范围与信息来源};

  % ---- 四种花法 ----
  \node[epi]  (g1) at (3.4, 3.2) {历史格 1\\重复运行，报告平均};
  \node[epi]  (g2) at (3.4, 1.1) {历史格 2\\预先给定，此后只读};
  \node[task] (g3) at (3.4,-1.1) {历史格 3\\任务内适配，用完即弃};
  \node[ver]  (g4) at (3.4,-3.2) {历史格 4\\跨任务持久记忆};

  \draw[arrow] (b) -- (g1);
  \draw[arrow] (b) -- (g2);
  \draw[arrow] (b) -- (g3);
  \draw[arrow] (b) -- (g4);

  % ---- 收益留到哪 ----
  \node[drop] (d1) at (6.8, 3.2) {不新增跨题状态};
  \node[drop] (d2) at (6.8, 1.1) {保留固定记忆};
  \node[drop] (d3) at (6.8,-1.1) {当前任务状态};
  \node[keep] (d4) at (6.8,-3.2) {保留到后续版本};

  \draw[arrow] (g1) -- (d1);
  \draw[arrow] (g2) -- (d2);
  \draw[arrow] (g3) -- (d3);
  \draw[arrow] (g4) -- (d4);

  % ---- 判据 ----
  \node[lbl, text width=37mm] at (10.5, 3.2) {提供静态参照；\\尚无重试候选选择};
  \node[lbl, text width=37mm] at (10.5, 1.1) {检查固定信息作用；\\需核对是否同源};
  \node[lbl, text width=37mm] at (10.5,-1.1) {检查任务内适配；\\不能据此声称跨题收益};
  \node[lbl, text width=37mm] at (10.5,-3.2) {检查持久方案；\\是否有收益仍需比较};

  % ---- 底部 ----
  \draw[gray!45] (-1.5,-4.6) -- (12.5,-4.6);
  \node[lbl, text width=140mm] at (5.5,-5.3) {
    设计要求：\textbf{（一）}预算单位与全流程成本范围明确；
    \textbf{（二）}各臂依赖的准备成本均入账；
    \textbf{（三）}任务内状态与跨题持久物分开；
    \textbf{（四）}开发改善不替代独立最终评价。分类图本身不证明预算对齐或收益。};
\end{tikzpicture}

\caption{历史E01四格的状态保留范围。静态重复取均值不同于生成多个候选后选择；固定记忆也可能帮助未见任务。四格分类不自动证明持久更新优越，完整预算与共同起点仍须核对。}
\label{fig:ch03-budget}
\end{figure}


\begin{table}[htbp]
\centering
\caption{历史E01四格的状态范围；机制分类不能代替收益检验}
\label{tab:ch03-fourarms}
\begin{small}
\begin{tabular}{>{\raggedright\arraybackslash}p{1.4cm}>{\raggedright\arraybackslash}p{3.4cm}>{\raggedright\arraybackslash}p{2.6cm}>{\raggedright\arraybackslash}p{4.4cm}>{\raggedright\arraybackslash}p{1.8cm}}
\toprule
格 & 执行方式 & 保留状态 & 比较用途与限制 & 本书实现 \\
\midrule
格 1 & 静态执行，重复取均值 & 不新增跨题记忆 & 提供静态成绩；不等于重试后选最佳 & \texttt{grid1} \\
\midrule
格 2 & 使用预先给定的只读持久物 & 固定记忆 & 检查已有信息的作用；历史起点未与格4同源 & \texttt{grid2} \\
\midrule
格 3 & 每道题单独适配，用完即弃 & 当前任务 & 区分任务内适配与跨题保留 & \texttt{grid3} \\
\midrule
格 4 & 从历史任务形成可保留记忆 & 后续系统版本 & 比较完整方案；收益与归因另行检验 & \texttt{grid4} \\
\bottomrule
\end{tabular}
\end{small}
\end{table}

\begin{sloppypar}
三条使用纪律：
\textbf{（一）}先固定预算单位和范围，并同时报告实际消耗；
\textbf{（二）}格2与格4各自依赖的准备过程都要入账；隐藏某一臂的准备成本，就会使该臂获得额外资源而不在账面体现；
\textbf{（三）}训练与选择集上的改善可以用于开发诊断，但泛化主张需要独立最终评价；不能将开发改善直接称为留出收益。
\end{sloppypar}

\begin{sloppypar}
外部harness重评论文比较了并行采样、顺序精化、harness演化与harness扩展，但这四条路径不是本书的四格\cite{harness-eval-2607}。其Table 1、2约束rollout次数与生成长度，不代表所有实际费用相同。另一个独立问题在Table 3中检验：先在训练题演化，再选版本并用于留出题，平均增益为\cited{+0.6}个百分点，其中GPT-5.4为\cited{+0.0}。这不是把同一演化产物换给新模型读取的实验，该表也没有并行采样臂。它提示跨题保留的修改仍可能缺乏泛化，不能据此替本书完成重试对照或跨模型迁移判决。
\end{sloppypar}

% ---------- 8. 实测 ----------
\section{本书的实测：E01 四格}
\label{sec:ch03-e01}

\begin{sloppypar}
历史E01四格运行的标识为 \texttt{E01-20260923T152230Z-90acda}。
条件：模型 \texttt{deepseek-flash}，解码温度 \measured{0.7}，
训练集 episode 配额 $E=$\measured{8}（四格同值），episode 内上限调用 \measured{8} / 步数 \measured{12}，
封存题 \measured{16} 道、每题重复运行 $R=$\measured{5}，重采样单位为题，协议所用 MDE 设计值 \calcnum{0.03}。
\end{sloppypar}

\textbf{训练划分的容量并不等于实际覆盖。}v1的训练划分有\measured{16}道题，但该次运行只取固定目录排序的前\measured{8}道：\texttt{classstate}、\texttt{dictmiss}、\texttt{errorcontract}、\texttt{mutabledefault}这\measured{4}个族，四格使用同一顺序。其余任务族没有进入这次训练。格4保存了\measured{8}条非空经验记录，来源回合判定为\measured{6}次通过、\measured{2}次未通过；旧提取器截取文本后逐条追加，没有独立的质量准入。失败回合也可能产生有用经验，成功回合也可能被错误概括，条目数不能直接当作有效经验数。后面的分族读数因此是在训练覆盖不均衡的条件下得到的；这一记录说明比较条件，不能单独解释某族分数为何升降。


\begin{table}[htbp]
\centering
\caption{E01四格主结果与账目。历史标签pass@1指一个Agent回合的通过率，回合内可有多次模型调用。通过数只计80个终评回合；主账回合计时累加8个训练与80个终评回合的记录值；每回合在最后一次验收前取值，也不含任务外适配、写入调用的等待时间。}
\label{tab:ch03-e01}
\begin{small}
\begin{tabular}{>{\raggedright\arraybackslash}p{1.2cm}>{\raggedright\arraybackslash}p{3.6cm}ccc}
\toprule
格 & 定义 & pass@1 & 通过 / 终评回合 & 回合计时和（秒） \\
\midrule
格 1 & 静态执行，重复取均值 & \measured{0.7375} & \measured{59}/\measured{80} & \measured{369.6} \\
格 2 & 持久物外部构建，只读 & \measured{0.7250} & \measured{58}/\measured{80} & \measured{417.7} \\
格 3 & 每道题单独适配，用完即弃 & \measured{0.7625} & \measured{61}/\measured{80} & \measured{411.3} \\
格 4 & 跨任务持久更新 & \measured{0.7875} & \measured{63}/\measured{80} & \measured{407.1} \\
\bottomrule
\end{tabular}
\end{small}

\vspace{2mm}
\begin{footnotesize}
数字身份：[实测] 本书运行 \texttt{E01-20260923T152230Z-90acda}；原始记录
\texttt{experiments/runs/E01-20260923T152230Z-90acda/}；成本账另见附录 \ref{app:cost}。
\end{footnotesize}
\end{table}

\begin{table}[htbp]
\centering
\caption{E01 判决输入：配对差的点估计与 95\% 自助区间（单位 = 题）}
\label{tab:ch03-verdict}
\begin{tabular}{>{\raggedright\arraybackslash}p{2.2cm}>{\raggedright\arraybackslash}p{1.6cm}>{\raggedright\arraybackslash}p{2.4cm}>{\raggedright\arraybackslash}p{4.6cm}>{\raggedright\arraybackslash}p{1.8cm}}
\toprule
比较 & 关联问题 & 点估计 & 95\% 区间 & 解读范围 \\
\midrule
格 4 − 格 1 & P1 & \measured{+0.0500} & \measured{$[-0.1000, +0.1750]$} & 方向未定 \\
格 4 − 格 2 & P1b & \measured{+0.0625} & \measured{$[-0.1375, +0.2750]$} & 方向未定 \\
格 4 − 格 3 & — & \measured{+0.0250} & \measured{$[-0.1250, +0.1750]$} & 方向未定 \\
\bottomrule
\end{tabular}
\end{table}

\begin{sloppypar}
三组比较目前均未提供足够依据支持方向性结论。
历史精度估算为 \calcnum{353} 道与 \calcnum{836} 道，现有 \measured{16} 道；
它们属于下一轮的规模参考，不能称为本轮预注册样本量。
此外，预算范围、基线行为和经验来源影响收益归属，见下面的实现核查。
\end{sloppypar}

\begin{sloppypar}
分族读数是这轮实测里最该被记住的部分（表 \ref{tab:ch03-family}）：
八个族每族 \measured{2} 道配对题，点估计从 \measured{$-0.3000$} 到 \measured{+0.4000} 都有。
它说明两件事：一是\textbf{每族两道题不足以提供稳定的族级总体估计}，可作探索描述及后续规模规划的输入，不能把不稳定方差当成已知参数；
二是“总体为正”完全可能由一两个族撑起来——这正是为什么本章要求按族报告，
以及为什么第 5 章要用族留出（而不是随机留出）来测迁移。
\end{sloppypar}

\begin{table}[htbp]
\centering
\caption{E01 分族读数（格 4 − 格 1）：每族 \measured{2} 道配对题，仅作探索描述}
\label{tab:ch03-family}
\begin{tabular}{>{\raggedright\arraybackslash}p{3.0cm}c>{\raggedright\arraybackslash}p{3.0cm}c}
\toprule
族 & 点估计 & 族 & 点估计 \\
\midrule
\texttt{classstate} & \measured{$-0.1000$} & \texttt{mutabledefault} & \measured{+0.4000} \\
\texttt{dictmiss} & \measured{+0.1000} & \texttt{offbyone} & \measured{$-0.1000$} \\
\texttt{errorcontract} & \measured{$0.0000$} & \texttt{recursionbase} & \measured{$-0.3000$} \\
\texttt{sortkey} & \measured{+0.3000} & \texttt{strindex} & \measured{+0.1000} \\
\bottomrule
\end{tabular}
\end{table}

\subsection{结果表中的名字必须与实际执行一致}

格 1 在旧标签中被写成“独立重跑，取最好的一次”，实际结果表却不是按最好一次计分。
核对该次运行记录的提交 \texttt{346742dc} 可见，
\texttt{run\_e01.py} 的 \texttt{sealed\_phase} 保存每题所有重复运行的布尔判定，
随后取平均通过率；没有跨重复运行挑选最佳答案的步骤。
因此 \measured{59}/\measured{80} 应解释为回合通过率的读数，不能解释为 best-of-$R$ 成功率。

本章据此把格 1 的实测标签校正为“静态执行器，重复运行取均值”。
它在一次回合内允许多步行动和执行测试，但没有单独实现一个把剩余总预算全部用于多候选重试的完整基线。
原蓝图中的预算四格仍是设计问题；当前实现只是其中一种受限操作化。
P1 的完整确认性比较还需明确重试次数、候选选择器、选择反馈以及总预算，不能只把格名改好就视为已补实验。

P1b 也存在类似的归属限制。
格 2 的外部经验与格 4 的学习经验并非同一初始工件，格 4 减格 2 不能纯粹隔离“写入开关”。
下一轮可从相同起始记忆分别运行冻结分支与更新分支，
同时报告来源构建成本与新经验质量；若仍使用不同来源，就应承认比较的是完整方案。
这些限制属于实验设计，增加重复次数不能消除它们。

\subsection{失败调用、缺失观测与零分分开处理}

E03c 曾有一个实验臂因账户返回 HTTP 402 而中断，后来已按留痕续跑完成；第5章保留这段历史及限制。这里用当时的缺失说明记账原则：账户拒绝不是 Agent 在该题上作答失败。
模型没有获得正常执行机会，因而不存在可直接填入通过率分母的完成回合。
把未运行全部填零会惩罚某个臂的账户状态；静默删除它们又会让读者误以为比较完整。
正确的报告应分别列出计划、启动、完成、缺失和重试数量，并保留基础设施错误。

一次已完成但答案错误的运行应进入任务成绩，网络拒绝或环境未启动则属于另一类记录。
两类记录都进入可用性与资源报告，但是否进入某个能力指标，须由预先声明的估计目标决定。
若目标是端到端服务可靠性，基础设施故障可以按事先规则构成失败；
若目标是模型更新机制，排除后的结果必须同时披露可用性，不能假装排除没有改变适用范围。

配对比较先核查每臂任务标识唯一、同一任务的族与划分一致，以及两侧都有非空的二值运行记录。当前\texttt{pair\_up}在这些输入检查后双向报告单侧缺项，再取共同任务；缺失如何计分仍由估计目标决定。
只分析完整配对可以避免捏造观测，却可能把最困难、最容易缺失的题排除出去。
报告除了差值，还应列出缺失发生在哪些族、由什么原因造成，以及完整配对覆盖了原计划的多少。
缺失与难度有关时，完整配对的成绩只能代表留下来的那部分任务。

续跑至少记录父运行、恢复时间、模型请求名、能取得的响应版本、代码和任务哈希、复用记忆哈希与新增回合。
跨日期使用相同模型别名，不能保证服务方没有更新后端。
因此补跑可以让试验矩阵完整，仍不能消除日期与实验臂重合的混淆；
确认性比较应尽量交错执行各臂，或设置同期基线来检查服务变化。

\subsection{从试验读数走向下一轮可检验设计}

E01 的下一步不是直接把现有模板复制到某个整数题数。
先定义推广范围，再核对题目是否独立、四格是否实现了对应干预，最后才计算规模。
本书下一轮应冻结一份能够交给其他人执行的计划，
并把已经看过的 v1 数据定位为试验与设计依据，不能再次称为新封存数据。

计划首先列出目标分布和共同起点。
P1 需要说明哪些预算用于更新、哪些用于重试，以及选择答案是否访问测试反馈；
P1b 需要说明只读分支与更新分支是否共享初始记忆。
若这些条件不同，就把相应比较标为完整方案比较，不能事后缩写成单机制效应。

计划随后固定主要比较、统计单位、效应门槛、区间方法和停止规则。
同一批数据支持多个命题时，需要说明哪些比较是主要检验，哪些只是探索；
如果要同时声称多条排序成立，应控制联合错误风险，不能拿每一项单独的区间直接拼成完整排序。
对于任务族很少的情形，应先增加独立族覆盖，而不是靠更多自助抽样次数获得看似精确的端点。

计划最后写清运行失败之后怎么办。
允许的重试条件、重试次数、恢复工件和预算上限在启动前固定，
所有尝试都进入原始账本；达到资源上限而未完成时，报告部分运行及其缺口。
这样安排得到的不一定是正结果，但会得到一个读者知道如何质疑、如何重复、如何继续的实验。


% ---------- 9. 条件变化 ----------
\section{固定版本的测试，与持续学习流程的测试}
\label{sec:ch03-streams}

前面的配对比较默认我们能够说清每个待比较方案。但在持续学习中，方案并不总是一个固定工件。同一条经验进入记忆以后，会影响后续检索、行动和新的经验来源；后面的结果又反过来影响下一次更新。因此，固定某个最终版本去做题，与让一个更新流程沿任务流继续运行，回答的是不同问题。

\subsection{先确定比较的是工件还是更新流程}

固定工件比较可以在开发结束后冻结两个版本，分别在留出任务上运行。每个任务开始时加载同一份组件，任务结果不再改变后面任务使用的版本。这种设计比较的是已经得到的两个对象。它适合回答候选是否值得发布，却不能直接说明生成该候选的整个流程是否更省资源、更快适应，或者在另一条任务顺序上仍然有效。

持续流程比较则固定更新规则，让组件状态按规则变化。可以把规则记作 $\pi$，把它在前 $i-1$ 个任务之后形成的状态记作 $A^{\pi}_{i-1}$。对于另一条更新规则 $\rho$，一个教学性的比较对象为
\begin{equation}
 \Delta_n^{\mathrm{flow}}=
 \mathbb E\!\left[\frac{1}{n}\sum_{i=1}^{n}
 \left\{q(A^{\pi}_{i-1},x_i)-q(A^{\rho}_{i-1},x_i)\right\}\right].
 \label{eq:ch03-flow-estimand}
\end{equation}
期望覆盖预先声明的任务流与运行随机性；预算规则也属于流程定义，只为简洁而没有写进每一项。这个记号说明比较对象，不提供任何性能保证。两条规则可以从相同初始组件出发，却因反馈和动作不同逐渐形成不同状态。到后半段时，不能再把它们解释成只有当前这一条记忆不同的局部比较。

式中的任务流若来自一个随时间变化的过程，平均值也会依赖观察时长和顺序。一个很快适应却后来退化的流程，与一个开始慢却稳定的流程，可能具有相同全程平均分。除总平均外，预先安排分阶段读数可以帮助识别这种差异；看完曲线以后再挑最有利区段，只能作为探索性分析，不能冒充原先要检验的主结果。

\subsection{把当前任务的评分与后续更新分开}

如果研究问题是“此前积累的经验能否帮助新任务”，当前任务的首次结果应在吸收该任务反馈之前确定。完成评分后，反馈才按协议进入下一状态。否则，系统先看当前题的验收信息再重新作答，报告的就是获得本题反馈后的适配能力。两者都可以研究，但不能把后者记作未经反馈的首次迁移表现。

任务内允许的工具检查需要单独说明。有些编程任务允许执行器运行公开测试，这属于当前任务的执行环境；最终评价还可以另设隐藏检查。系统在允许的公开测试中调试，不等于已经偷看最终答案。相反，把最终测试反馈通过另一条日志路径传回修改器，即使文件名仍叫“封存”，也已经改变了信息边界。应按实际可见信息定义任务条件。

若最终验收会影响后续更新，整段任务流就是在线评价与学习的组合设计，不能同时声称整段中的每一道题永远未被学习过程使用。可以明确规定：每题在被用于更新之前先记录一次预测或执行结果，此后允许其反馈影响未来任务。这样保留的是面向未来任务的时间边界，与开发完以后一次性打开的封存测试集不同。

反馈迟到时也不能简单按任务编号重新排序。某条流程在后续任务到来时尚未得到前一条反馈，另一条流程可能因执行较快已经得到。如果目标是比较真实服务流程，这种反馈可用时间可以是效果的一部分；如果只想比较更新机制，应另设计相同反馈释放时刻的受控比较。报告哪一种，取决于研究问题，不能在结果出现后随意切换。

\subsection{重复单位可能是一整条任务流}

在持续流程里，后一个任务的状态包含前面任务的影响。把一条长轨迹中的所有任务当作互不相关的重复样本，会丢掉这条依赖。增加同一条流的长度可以观察更长运行，却不等于获得更多独立的系统重启。独立重复应根据设计说明：是否重新初始化组件、如何生成任务流、哪些外部条件能够独立变化。

一种可审查设计是给两条规则共享任务流，在各自独立的状态副本中运行，然后对整条流的指标形成配对差。再用多条按预定方式产生的任务流，观察这种差在重复之间如何变化。共享输入不表示共享新生成的记忆；后者会使一个处理组泄漏信息给另一个处理组。任务生成与状态更新应有独立记录。

把同一组题换几个顺序，是研究顺序敏感性的一种办法，却没有增加题目内容的覆盖。可以报告“在这组题的不同顺序上怎样变化”，但不能把排列数量当作独立新任务族数量。若任务流来自同一个真实项目，项目内的相关性也不能因任务编号不同而消失。第\ref{ch13}章比较触发策略时尤其需要保留这些层次。

共享远程服务还可能使两组结果同时受到时段影响。把一个方案全部安排在早上、另一个全部安排在晚上，会混入时段条件；交错顺序可以减少某些系统性偏差，但不能保证服务完全稳定。固定模型版本、记录请求时间和说明无法控制的环境因素仍然必要。随机化与日志是互补手段，不能互相替代。

\subsection{中断会改变哪一项观测}

一条持续流程中途停止，至少留下三个不同事实：停止前已经完成的任务、没有完成的预定任务、以及形成到哪一步的组件状态。把未完成任务统一补零，可以描述某种“在给定时限内完成服务”的指标，但未必适合估计机制在已执行任务上的质量；把未完成任务全删掉，则会偏向容易早完成的轨迹。缺失处理应由原先的目标决定。

因此需要先记录中断原因。额度耗尽、网络故障、候选构造失败与主动达到停止条件，其解释不同。即使最后采用相同的计分规则，也应在报告中分列数量和原因，使读者能够判断结果是否被装置故障主导。第\ref{ch08}章E05的小规模试跑就说明了这个问题：更少的调用与更早的失败可以同时出现，不能仅凭资源账判断更有效率。

恢复以后，任务流究竟从哪个状态继续，需要由第\ref{ch02}章的版本和请求记录确定。重新从初始状态跑完整流程会产生一条新尝试；从检查点继续则需要保留之前的成本与可见反馈。两种做法都不能无说明地与中断前的记录拼成一个没有中断的样本。完整保留尝试史，是理解长期流程表现的前提。

\section{同一组差值，为什么会得到两个平均数}
\label{sec:ch03-weighting-example}

任务分布不只涉及从哪里取题，还涉及每种题在结论中占多大权重。下面构造一个只用于算术说明的例子，所有读数均为示例输入，不是模型运行结果。设两个任务族中，甲族有三道题，每题的方案差为\calcnum{+0.2}；乙族只有一道题，方案差为\calcnum{-0.4}。这些差值可以理解为题内重复运行均值之差，但本例没有实际进行重复运行。

若每道题权重相等，题均差为
\begin{equation}
 \bar d_{\mathrm{task}}=
 \frac{0.2+0.2+0.2-0.4}{4}=\calcnum{0.05}.
 \label{eq:ch03-task-weighted}
\end{equation}
若先在族内平均，再给两个族相同权重，则得到
\begin{equation}
 \bar d_{\mathrm{family}}=
 \frac{0.2+(-0.4)}{2}=\calcnum{-0.10}.
 \label{eq:ch03-family-weighted}
\end{equation}
两个结果的正负相反，却没有哪一步算错。第一种平均让题数更多的甲族占更大权重；第二种平均给两族相同权重。它们代表不同的目标任务混合。若实际部署更接近第一种混合，第二个数不自动更公平；若抽题数量只是采集便利造成的，第一种混合也不自动代表真实部署。

这个例子不能用来挑选有利的平均方式。正确顺序是先根据目标分布规定权重，再计算结果。若目标分布本来不确定，可以报告几个事先说明的权重情景，观察结论是否敏感；不能只展示能产生正向差值的那个情景。权重变化属于研究问题变化，不能被包装成一种普遍更正确的统计技巧。

一般地，若族权重为 $w_f$，族内平均差为 $\bar d_f$，则总体差写成 $\sum_f w_f\bar d_f$，其中权重非负且和为一。这个式子不要求所有族具有相同数量的任务，也不自动解决族内估计不精确的问题。给某族很大的权重，同时只测该族极少任务，会使总体结论依赖不稳定的一小块证据。报告分族样本数与差值，能够把这种依赖暴露出来。

同一道理也适用于不同成本区间或不同用户群。一个整体均值上升，仍可能掩盖某个重要子集的退化。是否允许这种交换应由任务要求说明，不能只靠总体平均决定。若协议设定了不可接受的子集退化条件，就要在候选选择与结果报告中保留它，而不是最后把所有指标加权成一个数再宣布通过。

本例只有示意性的点值，因此没有推断区间，也没有任何命题判决。真实实验还需要根据抽样方式、重复层次和相关结构估计不确定性。计算得到\calcnum{0.05}，只完成了“按某种权重如何汇总”的步骤；它既不证明真实总体差为正，也不证明该权重代表所关心的部署场景。算术、抽样与目标分布是三个需要分别检查的环节。


\section{缺失结果能把结论推到哪里}
\label{sec:ch03-missing-bounds}

失败与缺失分开记录之后，还需要回答一个实际问题：没有完成的观测，会不会改变我们对方案差异的解释？一个有用的起点不是立即选择填补模型，而是先利用指标本身的取值范围，检查缺失项在最好与最坏情况下能产生多大影响。这个检查不需要假定缺失是随机的，但给出的结论也很有限。

\subsection{先保留预定分母}

设原计划比较 $N$ 道配对任务，每题两方案的得分均落在 $[0,1]$，因而差值 $d_i$ 落在 $[-1,1]$。其中已完成的任务集合为 $\mathcal O$，还有 $m$ 个配对差完全未知。若目标仍是这批预定任务的平均差，那么它在已知记录条件下只能位于
\begin{equation}
 \frac{\sum_{i\in\mathcal O}d_i-m}{N}
 \ \leq\ \bar d_{\mathrm{planned}}\ \leq\
 \frac{\sum_{i\in\mathcal O}d_i+m}{N}.
 \label{eq:ch03-missing-bounds}
\end{equation}
左端把每个未知差值取为允许范围的下端，右端取为上端。这是由取值范围直接推出的界，不是置信区间，也没有使用某个覆盖概率。它暂时把已完成的差值当作已知输入；若这些差值本身来自有限重复估计，还要另外处理它们的不确定性。

下面给出构造的算术示例，数值不是本书实验读数。计划比较五道题，四道已知差值依次为\calcnum{0.2}、\calcnum{0.1}、\calcnum{0}和\calcnum{-0.1}，另一道完全缺失。只平均已完成的四道，得到\calcnum{0.05}；保留原来的五道题分母，则完整平均可能落在\calcnum{$[-0.16,0.24]$}。一个正的完整案例均值，并没有排除预定任务集上的负平均差。

这个例子没有说明缺失题一定很差，也没有证明两方案真实相同。它只说明：在不增加其他假设的情况下，现有信息不足以锁定差异方向。如果报告只写“平均提升五个百分点”，读者便看不到分母变化带来的这层不确定性。把计划数、完成数与完成案例结果并列，能使这个缺口直接可见。

\subsection{利用已知的一侧收紧界}

缺失并不总是两侧都没有结果。如果第五道题的基线得分已知为\calcnum{0.6}，只有候选得分未知，那么该题差值不再可以任意落在 $[-1,1]$，而是落在\calcnum{$[-0.6,0.4]$}。把这段范围代入同一个分母，完整平均的范围收紧为\calcnum{$[-0.08,0.12]$}。两种范围都来自同一组已完成差值，区别是第二种利用了更多真实已知信息。

一般地，如果每个缺失差值已经有合法下界 $l_i$ 和上界 $u_i$，就分别把它们相加后并入分子。界的依据可以来自指标范围、已完成的另一臂，或者确实已经记录的部分重复；不能来自“这道题看起来应该能做对”。对后一种预测可以另作敏感性分析，但应清楚标出它引入了额外假设。

这也说明为什么不能只保留最后的平均值。若逐题结果还在，就能知道缺的是哪一臂、哪些重复已经完成，因而可能得到更有信息量的范围。只留下“已完成若干题，平均多少”，后续往往只能采用更宽的界。完整记录没有增加新的模型能力证据，却避免了已有证据在汇总时被丢掉。

\subsection{界很宽时，下一步不一定是补跑所有题}

如果缺失范围已经覆盖会改变工程决策的结果，首先检查缺失原因和恢复可能性。原始结果已经生成但汇总失败时，可以修复汇总并核对原记录；任务根本未执行时，恢复汇总不会产生答案；服务返回格式错误时，则需要区分装置兼容性与任务能力。不同原因决定能否在不改变实验设计的前提下补齐观测。

补跑还要面对版本和时间条件。若原模型后端已经变化，补齐的新结果可能不再与原批次处于同一条件。保存它们仍然有价值，但需要报告续跑边界，必要时将其作为新的批次分析。不能因为原分母还差一项，就把任何后来的结果填入那个位置，假装条件从未变化。

如果无法补齐，可以继续报告范围与完成案例分析，并明确它们分别依赖哪些假设。对一个只关心有限任务集平均方向的问题，足够窄且完全同向的范围有时能给出有限结论；对推广到未来任务分布的问题，还缺抽样依据和总体不确定性。这两层不能因缺失处理完成而合并。

本节的上下界尤其不应被用来替代命题判决。它解决的是“既定任务集里缺掉的结果最多能改变多少”，并没有检验任务集是否代表目标分布，也没有修复经验来源、预算或反馈条件的混淆。一个设计失配的实验，即使所有结果齐全，仍可能无法识别原命题。


\subsection{不要把范围界误读成置信区间}

范围界与抽样区间来自不同的不确定性。范围界允许缺失项在合法区间内任意取值，追问的是当前记录与哪些完整结果相容；抽样区间则依赖抽样或运行模型，描述用有限数据估计某个目标量的精度。本节的方括号只有前一种含义。把它标成“95\%区间”，就为没有建立的概率解释加上了一个看似熟悉的数字。

如果已完成结果是每题若干次随机运行的均值，它们本身仍有波动。直接把这些均值放入公式，只得到给定这些估计值时的敏感性范围，不能保证覆盖总体的平均效果。要同时处理运行波动与缺失，需要先指定统计模型或有相应保证的方法，再解释最终区间覆盖的目标。不能把几个单独的区间端点拼起来，就默认整体仍具有相同覆盖率。

相反，如果目标仅是复核已经发生的有限批次，已完成的二值结果就是这批次的已知事实，不必为它们虚构抽样误差；未知的是没有发生或没有取得的结果。此时可以报告该有限批次的逻辑范围，但必须保留“这批次”限定。向未来任务、其他运行种子或其他服务版本推广，仍会重新引入需要分析的不确定性。

读者因此应先问区间在回答什么，再看它是否跨零。跨零的范围界说明缺失值足以改变当前汇总方向；跨零的统计区间说明在采用的统计方法下方向尚不确定。二者可能同时发生，也可能只有一种发生。明确这一区别，才能决定下一步更需要找回缺失记录、增加独立任务，还是修正比较设计。

\subsection{什么时候可以把未完成记为失败}

如果评价目标事先定义为“在规定时间内向使用者交付可用结果”，超时没有交付本来就可能属于失败事件。此时将其计为未成功，是执行目标的定义，而不是填补一个假装已知的模型答案。报告仍应保留超时原因，以便把服务可靠性与解题质量拆开观察。

如果目标是估计某种机制在完成执行时的任务表现，直接把网络中断记成程序错误，就混合了不同问题。只保留完成执行的任务也不自动解决，因为完成与否可能与题目难度或方案成本相关。更稳妥的报告方式是同时给出完成覆盖、失败原因和条件表现，并避免把条件表现外推成整个服务的成功率。

两种目标都可以成为研究问题，不能靠哪个数字更好看来选择。协议需要在运行前说明主结果如何处理超时、装置失败和用户撤销，分析时再如实列出偏离。若后来发现原定义不能服务工程决策，可以新增分析口径，但应保留原口径并说明更改原因，不能无痕替换。

\section{结果卡：让读者还原分母与比较对象}
\label{sec:ch03-result-card}

一张可读的结果表应让读者知道数字是什么，却不必把全部轨迹挤在正文里。可以在主表旁附上一段结果卡，按同样顺序交代比较对象、分析单位、实际覆盖、资源范围和证据限制。它不是用更多字段装饰表格，而是防止一个裸露的小数被不同读者理解成不同问题的答案。

比较对象首先写明两个完整方案。对E01，格4减格2涉及不同来源的经验库，不能只写“更新减不更新”；对E05，动态修订组是否真正获得过有效技能，需要从准入与检索记录判断，不能只看组名。方案名称负责定位实现，不能代替实际发生的处理。

分母要分层说明。任务族数、独立题数、每题重复数、候选数量和最终完成的episode数，不是可以互相替换的样本量。把它们合成一个很大的总数会隐藏结构；把所有成本除以这个数也可能得到无意义的平均。读者应能从结果卡回到逐题或逐流程记录，复算正文采用的那个分母。

主结果与诊断结果也应分开。首次成功率可以是主结果，格式错误次数可以是解释失败的诊断；修复后的成功率回答另一问题，需要显示额外调用。若一次试跑主要发现协议错误，最有价值的交付可能就是错误轨迹与修复验证，而不是强行从极小样本中提炼一种方法排序。

资源行应写计量单位与范围。调用次数、输入与输出token、墙钟和货币不能混写为一个没有单位的“成本”。未知账单就保持未知，失败候选仍进入相应成本范围。共享来源成本如何分摊，要能与研究问题对应；同一笔成本不能因为既出现在构建摘要又出现在总账中而重复相加。

最后一行写出当前结论到哪里为止，以及下一条什么证据可能改变它。例如“只验证单族运行路径”“直接复用失败尚未出现”“缺失范围跨过零”“模型别名未固定”，都比笼统的“有待进一步研究”更可执行。结果卡不必许诺下一项实验一定会产生支持性结论；它只需要让下一步能够针对当前缺口。


\section{条件变化：协议在什么情况下要改}
\label{sec:ch03-conditions}

\begin{sloppypar}
协议本身也有适用条件。三种变化最常见：
\end{sloppypar}

\begin{itemize}
  \item \begin{sloppypar}\textbf{换主干}。必须同时跑该主干的静态基线，并且把它写进同一张表
        （第 5 章 E03c 的 \texttt{base\_cross} 就是为此存在的）。
        只在新模型上跑带记忆的臂，得到的差里同时含“换了模型”与“有记忆”两件事。\end{sloppypar}
  \item \textbf{换任务族}。若要研究跨族迁移，就需按族留出；族内新题仍可研究族内泛化，但不能冒称跨族；
        代价是可用题数变少，样本量缺口更大。
  \item \textbf{换预算}。四格的 episode 配额与主账上限都要重新对齐；
        旧结果仍描述旧预算条件，不能改标签当作新预算结果。只有保存了足够轨迹且截断规则不改变此前行为时，才可做明确标注的离线截断分析；否则需要新运行。
\end{itemize}

\begin{pitfallbox}
\textbf{反例：把“最好的一次”写成“平均水平”。}
E01 格 1 的旧名称写“取最好”，实际计分却是重复回合取均值，表中据此命名。
若未来增加选优机制，必须连同选择器与额外预算一起单列比较，不能沿用现有读数。
同理，把 \measured{5} 次运行里最好的那一次挑出来当代表值，
等于把一个高方差量的\textbf{最大值}当成它的\textbf{均值}——
在本任务集上，\measured{16} 道题里有 \measured{9} 到 \measured{12} 道会随运行翻转，
挑选运行可能改变报告的排序与差异大小。
\end{pitfallbox}

\begin{pitfallbox}
\textbf{反例：把选择验证集上的提升充当收益。}
选择验证集是用来挑候选的，它的查询次数应受预定规则约束；达到额度后停止或按协议更换数据，不能只清零计数继续使用原题。
把在这一批题上观察到的上升当作最终收益，等于把搜索过程中最大的那一次抽样写成期望值。
训练与选择验证集上的改善也可以报告，但应标明开发用途；独立泛化主张需要相应留出评价，不能把开发收益换个名称当作封存结果。
\end{pitfallbox}

% ---------- 10. 能说与不能说 ----------
\section{本章能说什么，不能说什么}
\label{sec:ch03-claims}

\begin{sloppypar}
\textbf{能说：}
\end{sloppypar}
\begin{itemize}
  \item 在本任务集、本预算、本模型下，四格的封存集读数如表 \ref{tab:ch03-e01}，排序为格 4 最高。
  \item 三个配对差的 95\% 区间都跨过 0，\textbf{这些读数不足以裁定相应命题}。
  \item 历史区间精度估算为 \calcnum{353} 与 \calcnum{836} 道；它们不是完整功效设计。
  \item 观察到单题结果波动，估计成功概率需要重复：\measured{16} 道题里有 \measured{9} 到 \measured{12} 道随运行翻转。
\end{itemize}

\begin{sloppypar}
\textbf{不能说：}
\end{sloppypar}
\begin{itemize}
  \item 不能说"跨任务持久更新更好"——区间跨0，且重试策略、起始记忆与总预算的对照条件尚不齐全。
  \item 不能说“有持久物就有收益”——P1b需要同源起始记忆的只读/更新对照，该次试验尚不能分离该增量。
  \item 不能把本轮的排序外推到其他模型：E01未做换主干对照（那是第 5 章 E03c 的事）。
  \item 不能把 \measured{+0.0500} 这个点估计写进任何汇总表而不带它的区间与样本量。
\end{itemize}

\begin{principlebox}
自进化评测的第一条原则是：\textbf{先声明不判决的条件，再看数据}。
先跑数据、后找理由收紧或放宽门槛，得到的不是判决而是叙述。
下一轮应在启动前固定样本量规划、效应门槛、重采样单位与封存集口径，并保存当时的协议版本。
原运行已记录的设计字段与根据结果追加的规模计算要分开登记，后者不能追溯充当事前承诺。
\end{principlebox}

\falsifier{本章协议自身的失效条件：若某一组控制项无法在落盘文件里被指认
（模型、执行环境、任务分布、预算四项有任何一项只存在于论文描述里），
则不能声称该项混淆已被控制。仍可并列报告并计算描述性差值，但应标明条件缺失，不能把差值当成已识别的单机制效应。}

\openquestion{下一轮应先修复重试基线、共同起始经验与总预算控制，
再依据目标分布确定任务族、题数及重复次数。
需要补足的是可识别的比较和足够信息，不能只追求达到某个题数或让区间第一次不跨零。}

\section*{习题}
\begin{enumerate}[label=3-\arabic*.,ref=3-\arabic*]
  \item\label{ex:ch03-1} 表~\ref{tab:ch03-e01} 里格 4 的封存集读数比格 1 高 \measured{+0.0500}，
        区间跨过 0。列出三种与“跨任务持久更新”无关、却能产生同样差值的解释，
        并说明每种解释要补哪一个对照才能被排除。
  \item\label{ex:ch03-2} 设配对差标准差 $s=0.29$、目标区间半宽 $h=0.03$，
        按式 $n\ge(1.96\,s/h)^{2}$ 估算题数。若另要求双侧显著性水平 $0.05$、
        功效 $0.80$ 并检出真实差异 $0.03$，近似规模怎样变化？说明两个式子各自的假设。
  \item\label{ex:ch03-3} 一个团队报告"加入持久记忆后通过率从 \texttt{0.60} 升到 \texttt{0.71}"，
        但候选是在同一批题上挑出来的。按本章的四类题集，这句话缺哪两件东西才成立？
        补齐它们各要付出什么代价？
  \item\label{ex:ch03-4} 为什么换主干的静态基线必须是“同一张表里的另一行”，
        而不能引用该模型在别处的分数？用\cite{anatomy-memory-2602}表4的可恢复格式偏差与答题分说明，并指出不能由它们推出哪种因果份额。
  \item\label{ex:ch03-5} 你的资源只够把题数翻一倍。加在族内（每族题变多）还是加新族（族数变多）？
        它们分别增加哪一层的信息，对表~\ref{tab:ch03-metrics}中各项指标的估计与推广范围各有什么帮助和限制？
\end{enumerate}

\section*{常见误解}
\begin{itemize}
  \item \textbf{“区间跨过 0 就等于没有效果”。}
        跨0说明当前区间兼容两种方向；原因可能涉及效应大小、方差、抽样结构和样本量。点估计仍在，
        \measured{+0.0500} 与 \measured{+0.0625} 没有被否定，被否定的是“现在就能下结论”。
  \item \textbf{“做了配对比较就可以少跑几次”。}
        配对可抵消部分共同题目波动，不能消除所有题间差异或运行波动：
        本任务集 \measured{16} 道题里有 \measured{9} 到 \measured{12} 道随运行翻转。
  \item \textbf{“封存集只要没被训练过就够了”。}
        还要查最终评测启动前，它是否曾被开发与选择流程查询。
        反复查询会使选择过程适应该数据，削弱其作为独立泛化评价的资格；它与训练题泄漏的路径不同，都需检查实际信息流。
  \item \textbf{“预算对等是默认成立的”。}
        调用次数、步数、token、墙钟要逐项报告，并明确主要控制哪一种资源；
        只对齐其中一项，仍可研究这一约束下的方案表现，但不能把它说成所有资源都相等的机制比较。
\end{itemize}

\section*{进一步阅读}
\begin{itemize}
  \item \cite{anatomy-memory-2602}：主干变化与可恢复格式偏差（表4）的关系，以及
        “维护成本没有计入”这句自陈（表 5 附近）。本书“跨模型先跑静态基线”与“异步不等于免费”
        两条纪律都从这两处来。
  \item \cite{harness-eval-2607}：区分rollout约束下的四条路径与独立留出测试，
        留出题平均增益\cited{+0.6}与GPT-5.4增益\cited{+0.0}的比较对象；不将后者称为换读取者迁移。
        读它时应检查每张表实际包含的比较路径，不能把外部四条路径直接替换成本书四格。
  \item 附录~\ref{app:stats}：配对差、自助区间与样本量反推的完整推导，
        本章只用了结论式。
  \item 附录~\ref{app:lab}~\S\ref{sec:appb-tasks}：任务集四类划分与
        \texttt{SPLIT\_BY\_INDEX} 的实现位置，以及 $Q$ 预算抛异常的字段在哪。
\end{itemize}
